% !PRINTED 514
\weberpage{514}
%<<CONT
schaft haben, durch einander theilbar, also associirt, und wir stellen also die Definition auf:

1. \gesperrt{Das Functional $\delta = \alpha x + \beta y + \cdots$ und jedes damit associirte Functional heisst grösster gemeinschaftlicher Theiler von $\alpha, \beta, \dots$}

Wenn $\delta$ eine Einheit ist, so sagen wir auch, $\alpha, \beta, \dots$ seien ohne gemeinsamen Theiler, denn dann giebt es ausser den Einheiten kein ganzes Functional, das in allen $\alpha, \beta, \dots$ aufgeht.

2. \gesperrt{Zwei Functionale $\alpha, \beta$, die keinen gemeinsamen Theiler haben, für die also $\alpha x + \beta y$ eine Einheit ist, heissen relative Primfunctionale oder theilerfremd.}

Aus diesen Definitionen ergeben sich sehr einfach folgende Sätze:

3. \gesperrt{Wenn ganze Functionale $\xi, \eta, \dots$ in $\Omega$ existiren, derart, dass
\[ \alpha \xi + \beta \eta + \cdots = \varepsilon \]
eine Einheit ist, so sind die ganzen Functionale $\alpha, \beta, \dots$ ohne gemeinsamen Theiler.}

Denn haben $\alpha, \beta, \dots$ einen gemeinsamen Theiler $\delta$, so ist (nach \S.\,138) $\delta$ auch Theiler von $\alpha \xi + \beta \eta + \cdots$, und $\delta$ muss also auch eine Einheit sein.

4. \gesperrt{Sind $\alpha, \beta, \gamma$ drei ganze Functionale und ist $\alpha$ relativ prim zu $\beta$ und zu $\gamma$, so ist $\alpha$ auch relativ prim zu $\beta \gamma$.}

Denn nach Voraussetzung sind, wenn $x, y, u, v$ vier Variable sind,
\[ \varepsilon = \alpha x + \beta y, \quad \varepsilon_1 = \alpha u + \gamma v \]
Einheiten. Demnach ist auch
\[ \alpha (\alpha u x + \gamma v x + \beta u y) + \beta \gamma v y = \varepsilon \varepsilon_1 \]
eine Einheit. Da aber $\alpha u x + \gamma v x + \beta u y$ und $v y$ ganz sind, so folgt nach dem Satze 3., dass $\alpha$ und $\beta \gamma$ relativ prim sind.

Hieran schliesst sich der Beweis des folgenden sehr wichtigen Satzes:

5. \gesperrt{Sind $\alpha, \beta, \mu$ ganze Functionale, $\alpha$ relativ prim zu $\beta$ und $\alpha \mu$ durch $\beta$ theilbar, so ist $\mu$ durch $\beta$ theilbar.}
